\begin{description}
\item[(T10.1)] The rays travelling closer to the gravitational body will bend
more. Thus, we get shortest convergence point where the rays just grazing the
solar surface will meet each other.
% FIGURE NEEDED: ray grazing the Sun (radius $R_\odot$) bent by $\theta_b$,
% crossing the axis at distance $f_{\min}$ (2016_theory_sol.pdf p. 14).
\begin{align*}
\theta_{\mathrm{b}} &= \frac{2R_{\mathrm{sch}}}{R_\odot}
  \simeq \frac{R_\odot}{f_{\min}} \tag*{(2.0)}\\
\therefore f_{\min} &= \frac{R_\odot^2}{2R_{\mathrm{sch}}}
  = \frac{R_\odot^2 c^2}{4GM_\odot} \tag*{(2.0)}\\
  &= \frac{(6.955 \times 10^{8} \times 2.998 \times 10^{8})^2}
          {4 \times 6.674 \times 10^{-11} \times 1.989 \times 10^{30}}~\mathrm{m}\\
  &= 8.188 \times 10^{13}~\mathrm{m}
   = \frac{8.188 \times 10^{13}}{1.496 \times 10^{11}}~\mathrm{AU}
\end{align*}
\[
\boxed{f_{\min} = 547.3~\mathrm{AU}} \tag*{(2.0)}
\]

\item[(T10.2)] The following figure needs to be drawn. \hfill(1.0)
% FIGURE NEEDED: two rays with impact parameters $R_\odot$ and $R_\odot + h$
% (annulus of width $h$), detector of radius $a$ at $f_{\min}$, outer ray
% crossing the axis at $f_2$; bending angle $\theta_2$ marked
% (2016_theory_sol.pdf p. 14).

The light bending from the surface of the Sun ($r = R_\odot$) will intersect
the detector at its centre, as it is kept at $f_{\min}$.

The boundary of the detector will be intersected by a light ray with
$r = R_\odot + h$. This ray will intersect the $x$-axis at a distance $f_2$.
\[
f_2 = \frac{(R_\odot + h)^2}{2R_{\mathrm{sch}}} \tag*{(2.0)}
\]
\textbf{Same argument as in Part 1.} For small angles,
\begin{align*}
a &= (f_2 - f_{\min})\theta_2\\
  &= \left[\frac{(R_\odot + h)^2}{2R_{\mathrm{sch}}}
     - \frac{R_\odot^2}{2R_{\mathrm{sch}}}\right]
     \frac{2R_{\mathrm{sch}}}{(R_\odot + h)}
   = \frac{2R_\odot h + h^2}{R_\odot + h}\\
  &\simeq 2h \tag*{(2.0)}
\end{align*}
Let original intensity of the incoming radiation be $I_0$.

The flux at the detector in the presence of Sun is
$\Phi_\odot = I_0\, 2\pi\, R_\odot\, h$ \hfill(1.0)

The flux at the detector in the absence of Sun is
$\Phi_0 = I_0\, \pi\, a^2$ \hfill(1.0)

The amplification is therefore
\[
A_{\mathrm{m}} = \frac{\Phi_\odot}{\Phi_0}
  = \frac{I_0 2\pi R_\odot h}{I_0 \pi a^2}
  = \boxed{\frac{R_\odot}{a}} \tag*{(1.0)}
\]

\item[(T10.3)]
% FIGURE NEEDED: spherical mass distribution with two parallel rays at impact
% parameters $r_1$, $r_2$ bent by $\theta_1$, $\theta_2$, both converging at
% the same focus, $f_1 = f_2$ (2016_theory_sol.pdf p. 15).
All rays should focus at the same spot. This should be evident from figure
drawn on answersheet or otherwise. \hfill(2.0)

Let there be two rays with impact parameters $r_1$ and $r_2$. The
corresponding distances of focus will be
\[
f_i = \frac{r_i^2}{2r_{\mathrm{sch}_i}} = \frac{r_i^2 c^2}{4GM(r_i)} \tag*{(2.0)}
\]
\textbf{Same argument as in Part 1.} The rquirement % sic: "rquirement" in source
$f_1 = f_2$ implies
\[
\frac{r_1^2}{r_2^2} = \frac{M(r_1)}{M(r_2)} \tag*{(1.0)}
\]
The required mass distribution is: $\boxed{M(r) \propto r^2}$ \hfill(1.0)
\end{description}
